\clearpage
\begingroup
\fancyhead[L]{\small POSTHUMOUS PAPERS AND CORRESPONDENCE}
\renewcommand{\headrulewidth}{0pt}
\par\medskip
\Needspace{8\baselineskip}
\pdfbookmark[2]{[2] Nicolai to Gauss, 19 March 1814}{pallas-VIII-letter2}
\begin{center}
[2.]\par\medskip
\textsc{Nicolai} to \textsc{Gauss}, Seeberg, 19 March 1814.
\end{center}
\par\noindent\hspace*{\parindent}\{.\ .\ .\ .\ .\quad I have recently taken up the perturbations of Pallas
due to Saturn again, and I am pleased to send you today the completed results for three
elements. You may be surprised at
the considerable number of equations, which surprised me at first as well.
Most are mere trifles that will probably make little difference to the results.
Nevertheless, I have neglected \textls[100]{not a single} equation that for \(\log a\) exceeds
one unit or for the other elements exceeds \(0\overset{\prime\prime}{,}1\), just as
you did in the case of Jupiter. I can also vouch for the correctness of
\textls[100]{every} equation, since I first calculated them all with great care and
diligence and then checked them once more briefly. I
felt obliged to do this, since with the possible exception of the interpolations
there is no check throughout the calculation.\quad .\ .\ .\ .\ .\quad I must also mention
that I have taken as Saturn's mass the value that
\textsc{Laplace} gives, on the basis of more recent investigations, in the latest, 4\textsuperscript{th} edition of his
Exposition du Système du monde, which arrived here a few weeks ago,
namely \(\frac{1}{3512{,}08}\mathpunct{\raisebox{.45ex}{.}}\) --- Calculation by this other method
\SourceLeaf{586}
seems to me, incidentally, far simpler and more convenient than calculation by
the first method. There is also undoubtedly far less risk of errors.
Only the second part of the epoch worries me a little. --- The secular changes
for inclination and the node, as well as the larger equations,
agree quite well with my first calculation. --- I should have thought
that for perihelion and for the 2\textsuperscript{nd} part of the epoch the assumed
number of twelve positions of Pallas taken as a basis would not be
too small. But when I compare the number of equations for these two elements
in my earlier calculation with the number for the other elements in the same
calculation, and then consider the number of equations for the
latter that the new calculation has yielded, I am almost afraid
it will be. For that reason I now wish I had taken at least 18 positions of Pallas
as a basis. The continuation of the calculation will clarify all this.
The numerical calculation of all the equations for the 9 oppositions
will, incidentally, take a little more work.\quad .\ .\ .\ .\ .
\par\medskip
\Needspace{28\baselineskip}
\begin{center}
\(\operatorname{Log}a\) and first part of the epoch\HistoricalDisplayNoteMark.
\end{center}
\HistoricalDisplayNoteText{The expression on the left gives the perturbations of \(\log a\) in units of the
7th decimal place. The coefficients on the right, multiplied by the
corresponding cosine, give the perturbations \(\delta\int n\,\mathrm{d}t\).]}
\begingroup
\small
\newcommand{\VIIIpallas}{\text{\AstroSymbols ⚴}}
\newcommand{\VIIIsaturn}{\text{\AstroSymbols ♄}}
\setlength{\arraycolsep}{2pt}
\renewcommand{\arraystretch}{1.25}
\[
\begin{array}{r@{\,}l@{\qquad\quad}r}
5&\sin(\VIIIpallas+79^0\ 15')&0\overset{\prime\prime}{,}39\\
+7&\sin(2\VIIIpallas+93\ 48)&0{,}24\\
+2&\sin(3\VIIIpallas+341\ 19)&0{,}04\\
+2&\sin(\VIIIpallas+\VIIIsaturn+87\ 31)&0{,}15\\
+2&\sin(\VIIIpallas-\VIIIsaturn+92\ 48)&0{,}20\\
+2&\sin(2\VIIIpallas-\VIIIsaturn+77\ 14)&0{,}07\\
+2&\sin(3\VIIIpallas-\VIIIsaturn+88\ 13)&0{,}05\\
+1&\sin(\VIIIpallas+2\VIIIsaturn+140\ 23)&0{,}06\\
+38&\sin(\VIIIpallas-2\VIIIsaturn+34\ 30)&3{,}92\\
+39&\sin(2\VIIIpallas-2\VIIIsaturn+88\ 29)&1{,}65\\
+9&\sin(3\VIIIpallas-2\VIIIsaturn+331\ 1)&0{,}23\\
+2&\sin(4\VIIIpallas-2\VIIIsaturn+190\ 0)&0{,}03\\
+10&\sin(\VIIIpallas-3\VIIIsaturn+84\ 44)&1{,}37\\
+6&\sin(2\VIIIpallas-3\VIIIsaturn+84\ 30)&0{,}27
\end{array}
\]
\SourceLeaf{587}
\[
\begin{array}{r@{\,}l@{\qquad\quad}r}
+9&\sin(3\VIIIpallas-3\VIIIsaturn+81^0\ 4')&0\overset{\prime\prime}{,}24\\
+3&\sin(4\VIIIpallas-3\VIIIsaturn+325\ 50)&0{,}06\\
+1&\sin(\VIIIpallas-4\VIIIsaturn+113\ 8)&0{,}28\\
+2&\sin(2\VIIIpallas-4\VIIIsaturn+81\ 57)&0{,}09\\
+2&\sin(3\VIIIpallas-4\VIIIsaturn+82\ 45)&0{,}05\\
+2&\sin(4\VIIIpallas-4\VIIIsaturn+72\ 40)&0{,}04.
\end{array}
\]
\Needspace{10\baselineskip}
\begin{center}
Inclination of the orbit.\par
Annual change \(=+0\overset{\prime\prime}{,}16950\)
\end{center}
\[
\begin{array}{r@{\,}l@{\qquad}r@{\,}l}
0\overset{\prime\prime}{,}42&\sin(\VIIIpallas+42^0\ 58')&+0\overset{\prime\prime}{,}11&\sin(\VIIIpallas+2\VIIIsaturn+129^0\ 41')\\
+0{,}24&\sin(2\VIIIpallas+98\ 47)&+0{,}56&\sin(\VIIIpallas-2\VIIIsaturn+243\ 21)\\
+1{,}02&\sin(\VIIIsaturn+37\ 58)&+0{,}29&\sin(2\VIIIpallas-2\VIIIsaturn+275\ 8)\\
+0{,}18&\sin(\VIIIpallas+\VIIIsaturn+89\ 3)&+0{,}23&\sin(3\VIIIsaturn+83\ 5)\\
+0{,}11&\sin(\VIIIpallas-\VIIIsaturn+355\ 47)&+0{,}52&\sin(\VIIIpallas-3\VIIIsaturn+260\ 41)\\
+2{,}27&\sin(2\VIIIsaturn+105\ 16)&+0{,}13&\sin(\VIIIpallas-4\VIIIsaturn+273\ 9).
\end{array}
\]
\Needspace{12\baselineskip}
\begin{center}
Longitude of the ascending node.\par
Annual change \(=-1\overset{\prime\prime}{,}47934\)
\end{center}
\[
\begin{array}{r@{\,}l@{\qquad}r@{\,}l}
0\overset{\prime\prime}{,}89&\sin(\VIIIpallas+281^0\ 36')&+1\overset{\prime\prime}{,}22&\sin(\VIIIpallas-2\VIIIsaturn+111^0\ 22')\\
+0{,}39&\sin(2\VIIIpallas+11\ 13)&+0{,}59&\sin(2\VIIIpallas-2\VIIIsaturn+184\ 51)\\
+1{,}79&\sin(\VIIIsaturn+246\ 49)&+0{,}31&\sin(3\VIIIsaturn+327\ 21)\\
+0{,}17&\sin(\VIIIpallas+\VIIIsaturn+350\ 51)&+0{,}36&\sin(\VIIIpallas-3\VIIIsaturn+3\ 31)\\
+0{,}75&\sin(\VIIIpallas-\VIIIsaturn+186\ 33)&+0{,}14&\sin(2\VIIIpallas-3\VIIIsaturn+143\ 10)\\
+0{,}13&\sin(2\VIIIpallas-\VIIIsaturn+301\ 7)&+0{,}11&\sin(3\VIIIpallas-3\VIIIsaturn+180\ 41)\\
+3{,}27&\sin(2\VIIIsaturn+2\ 51)&+0{,}13&\sin(\VIIIpallas-4\VIIIsaturn+25\ 2).\}\\
+0{,}15&\sin(\VIIIpallas+2\VIIIsaturn+37\ 19)&&
\end{array}
\]
\endgroup
\clearpage
\endgroup
